\section{Implementación computacional}

El proyecto se implementó en Python y PyTorch. La diferenciación automática permitió calcular derivadas temporales de los estados aproximados y gradientes de los parámetros entrenables \citep{Paszke2019}. El repositorio contiene notebooks reproducibles, código modular del modelo Balloon--Windkessel, métricas, pruebas estadísticas y pruebas unitarias.

\subsection{GLM}

El GLM sintético utilizó una HRF canónica tipo SPM. En datos reales, el diseño incluyó seis regresores de tarea, 12 regresores de movimiento, tres términos de deriva y una constante. Las métricas de generalización se calcularon sobre el bloque retenido.

\subsection{MLP}

La MLP utilizó como entrada una historia temporal del estímulo y dos capas ocultas con activación hiperbólica. Este modelo no contiene estados fisiológicos ni restricciones diferenciales, por lo que funciona como referencia puramente basada en datos.

\subsection{PINN}

La red de estados contiene tres capas ocultas de 64 neuronas con activación $\tanh$ y precisión de 64 bits. El entrenamiento se realizó en dos fases con Adam, seguido de L-BFGS. La primera fase priorizó la estabilización de los estados; la segunda incrementó el peso de los residuos físicos y optimizó conjuntamente los parámetros.

\begin{table}[H]
\centering
\caption{Configuración principal de la PINN utilizada en el lote real.}
\label{tab:pinn_config}
\begin{tabular}{lc}
\toprule
Elemento & Configuración \\
\midrule
Capas ocultas & $64\times64\times64$ \\
Activación & $\tanh$ \\
Precisión & float64 \\
Adam fase 1 & 1.200 épocas \\
Adam fase 2 & 3.500 épocas \\
L-BFGS & hasta 300 iteraciones \\
Puntos de colocación & 700 \\
Parámetros estimados & $\epsilon,\tau,\alpha$ \\
\bottomrule
\end{tabular}
\end{table}

\subsection{Reproducibilidad y control de versiones}

El repositorio público contiene 13 notebooks finales, archivos de configuración, código modular y pruebas automáticas. Los datos HCP crudos, modelos binarios y archivos volumétricos se excluyeron por tamaño y condiciones de uso. Dos pruebas unitarias verifican el equilibrio basal y la respuesta positiva frente a un estímulo; ambas se ejecutaron correctamente en Python 3.11.
